\documentclass{article}
\usepackage[T1]{fontenc}
\usepackage{lmodern}
\usepackage{graphicx}
\usepackage{amsmath,amssymb}
\usepackage[a4paper,margin=2cm]{geometry}
\usepackage[absolute,overlay]{textpos}
\setlength{\TPHorizModule}{1mm}
\setlength{\TPVertModule}{1mm}
\usepackage[indonesian]{babel}
\usepackage{hyperref}
\setlength{\emergencystretch}{2em}
% unit: Halaman 44

\begin{document}

% === Halaman 44 (python/native) ===

AES-256

• AES-256  dengan AES-128 hanya berbeda dalam hal:

    1. Jumlah putaran 14 kali (AES-128 hanya 10 kali)

    2. Panjang kunci 256 bit (AES-128 hanya 128 bit)

    3. Expansion key (untuk membangkitkan kunci putaran)

• Expansion key pada AES-128 uritannya adalah sebagai berikut:

       RotWord → SubWord → XOR Rcon → XOR dengan word sebelumnya

• Expansion key pada AES-256: 

      Operasi SubWord tambahan pada word tertentu:

         jika i mod 8 = 0, dilakukan RotWord → SubWord → XOR Rcon 

         jika i mod 8 = 4, dilakukan SubWord saja

\end{document}
